\usepackage[utf8]{inputenc}
\usepackage{amsmath}
\usepackage{ragged2e}
\usepackage{varwidth}
\usepackage{indentfirst}
\usepackage{calc}
\usepackage{tabularx}

\usepackage[small]{titlesec}
\usepackage[a4paper,
    left=30mm,
    right=15mm,
    top=20mm,
    bottom=20mm]{geometry}

\justifying
\sloppy

\usepackage{fontspec}

\defaultfontfeatures{Ligatures={TeX},Renderer=Basic}
\setmainfont[Ligatures={TeX,Historic}]{Times New Roman}
\setlength{\parindent}{1.25cm}

\linespread{1}

% hyperlinks
\usepackage[final,hidelinks]{hyperref}

\renewcommand{\UrlFont}{}

% resourses list
\usepackage[square,numbers,sort&compress]{natbib}

\bibliographystyle{belarus-specific-utf8gost780u}
\setlength{\bibsep}{0em}
\def\BibUrlRus#1{--- Режим доступа~: \url{#1}}
\def\BibDateRus#1{--- Дата доступа~: \parsedate{#1}}
\def\BibUrl#1{\BibUrlRus{#1}}
\def\BibDate#1{\BibDateRus{#1}}
\def\BibTitleFormat#1{\StrSubstitute{#1}{ - }{ --- }}

% page numbering
\usepackage{fancyhdr}

\pagestyle{fancy}
\fancyhf{}
\fancyfoot[R]{\thepage}
\renewcommand{\headrulewidth}{0pt}
\addtocontents{toc}{\protect\thispagestyle{fancy}}

% lists
\usepackage{enumitem}

\setlist{
    topsep=0em,
    itemsep=0em,
    parsep=0em,
    listparindent=\parindent
}

\setlist[itemize,0]{itemindent=\parindent + 2.2ex, leftmargin=0ex, label=--}
\setlist[enumerate,1]{itemindent=\parindent + 2.7ex, leftmargin=0ex, label=\arabic*}
\setlist[enumerate,2]{itemindent=\parindent + \parindent - 2.7ex, label=\arabic*}
\renewcommand{\labelitemi}{\textendash}

% code listing
\usepackage{listings}
\usepackage[dvipsnames]{xcolor}

\lstset{
    columns=fullflexible,
    basicstyle=\ttfamily\small,
    tabsize=4,
    breaklines=true,
    breakatwhitespace=false,
    showstringspaces=false,
    language=Java,
    linewidth=\textwidth,
    commentstyle=\color{gray},
    keywordstyle=\color{orange},
    stringstyle=\color{Green},
}

% sections and toc
\newcommand{\sectionuppercase}[1]{\section[#1]{\uppercase{#1}}}

\makeatletter
\renewcommand{\@pnumwidth}{1em}
\renewcommand{\@dotsep}{1}
\renewcommand{\l@section}{\@dottedtocline{1}{0.5em}{1.2em}}
\renewcommand{\l@subsection}{\@dottedtocline{2}{1.7em}{2.0em}}
\makeatother

\titleformat*{\section}{\bfseries\MakeUppercase}

% images
\usepackage{graphicx}
\usepackage{float}

\graphicspath{ {images/} }

% images and tables caption
\usepackage{caption}
\usepackage{subcaption}

% FIXME: more than one line spacing after caption

\setlength{\intextsep}{\baselineskip}

\DeclareCaptionLabelFormat{stbfigure}{Рисунок #2}
\DeclareCaptionLabelFormat{stbtable}{Таблица #2}
\DeclareCaptionLabelSeparator{stb}{~--~}
\captionsetup{labelsep=stb}
\captionsetup[figure]{
    labelformat=stbfigure,
    aboveskip=\baselineskip,
    belowskip=0pt,
    justification=centering
}

% for table use tabularx
\captionsetup[table]{
    labelformat=stbtable,
    aboveskip=0pt,
    format=hang,
    justification=raggedright,
    singlelinecheck=false
}
\renewcommand{\thesubfigure}{\asbuk{subfigure}}

\AtBeginDocument{\numberwithin{table}{section}}
\AtBeginDocument{\numberwithin{figure}{section}}

% sections spacing and styles
% TODO: rewrite

\newlength{\fivecharsapprox}
\setlength{\fivecharsapprox}{1.25cm}

\titlespacing{\sectioncentered}
{0em}
{-1.5em plus -1ex minus -.2ex}
{1em plus .2ex}

\titlespacing{\section}
{\fivecharsapprox}
{-1.5em plus -1ex minus -.2ex}
{1em plus .2ex}

\titlespacing{\subsection}
{\fivecharsapprox}
{1em plus 1ex minus .2ex}
{1em plus .2ex}

\titleformat{\subsubsection}[runin]
{\pretolerance=10000\filright\normalfont}
{\bfseries\thesubsubsection}
{0.5em}
{}
[\hspace*{0em}]

\titlespacing{\subsubsection}
{\fivecharsapprox}
{0em}
{0em}

% annex section
% TODO: rewrite

\usepackage{totcount}

\makeatletter
\def\@annexAlph#1{%
    \ifcase#1\or А\or Б\or В\or Г\or Д\or Е\or Ж\or И\or К\or Л\or%
        М\or Н\or П\or Р\or С\or Т\or У\or Ф\or Х\or Ц\or Ш\or Щ\or Э\or%
        Ю\or Я\else\@ctrerr\fi%
}
\def\annexAlph#1{\expandafter\@annexAlph\csname c@#1\endcsname}
\makeatother

\newtotcounter{equation_annex}

\newcommand{\addtocline}[2]{
    \phantomsection
    \addcontentsline{toc}{#1}{#2}
}

\newcommand{\annex}[2]{
    \newpage
    \stepcounter{equation_annex}
    \centerline{\textbf{ПРИЛОЖЕНИЕ \annexAlph{equation_annex}}}
    \centerline{(\textbf{#1})}
    \addtocline{section}{Приложение \annexAlph{equation_annex} (#1) #2}
    \centerline{\textbf{#2}}
}
